%%%%%
%
\section{Cr�ation d'un vecteur composite}

\itemstyle{Op\'erateur~:} \textsf{select\_sur\_vect}
\index{select\_sur\_vect}
\label{select_sur_vect}

%%%%%OBJET%%%%
\itemstyle{Objet~:}
Cr�ation d'un vecteur composite, qui contient les donn�es de $X$ 
si elles sont pr�sentes et comprises entre crit\_min 
et crit\_max, et les donn�es de $Y$ (sans filtrage) dans le cas contraire.

%%%%%ARGUMENTS%%%%
\bigskip
%%%%%ENTREES%%%%
\begin{tabarg}
\multicolumn{3}{l}{\titretab {Entr\'ees}}\\
\hline
$X$ & Vecteur, \texttt{float} & Param�tre en entree\\
$Y$ & Vecteur, \texttt{float} & Param�tre de remplacement (m�me unit� que $X$)\\
$Crit\_{min}$ & Coeff., \texttt{float} & Valeur mini pour laquelle $Z=X$\\
$Crit\_{max}$ & Coeff., \texttt{float} & Valeur maxi pour laquelle $Z=X$\\
\end{tabarg}

\medskip
%%%%%SORTIES%%%%
\begin{tabarg}
\multicolumn{3}{l}{\titretab {Sorties}}\\
\hline
$Z$ & Vecteur, \texttt{float} & S�rie de donn�es composites (m�me unit� que $X$)\\
\end{tabarg}


%%%%%FORMULE%%%%
\itemstyle{Formule}

%%select_sur_vect
1- Recherche des �ventuelles valeurs hors gamme dans le vecteur $X$ \\
2- Remplacement dans ce cas par des valeurs du vecteur $Y$ \\

Un crit�re peut �tre non renseign� : l'utilisateur doit alors
donner pour valeur (0.0/0.0).


%%%%%COMMENTAIRE%%%%
\itemstyle{Commentaire}

Ne pas oublier que les vecteurs $X$ et $Y$ doivent avoir le m�me 
nombre d'�chantillons.

Cet op�rateur peut �tre consid�r� comme une version plus puissante de
l'op�rateur \textsf{Complete\_X\_avec\_Y} (\ref{Complete_X_avec_Y}) : ce 
dernier ne traite que les valeurs manquantes, alors que 
\textsf{select\_sur\_vect} applique aussi un filtre simple (valeur
entre deux bornes). Mais il a un gros d�faut : il ne produit pas de 
flag indiquant l'origine de chaque valeur.

%%%%%AUTEUR%%%%
\itemstyle{Auteur~:}
CNRM/GMEI/TRAMM

%%%%%REFERENCE%%%%
\itemstyle{R\'ef\'erences~:}

